% Part: computability
% Chapter: computability-theory
% Section: halting-problem

\documentclass[../../../include/open-logic-section]{subfiles}

\begin{document}

\olfileid{cmp}{thy}{hlt}
\olsection{నిలుపు సమస్య (హాల్టింగ్ సమస్య)}

నిర్మాణం ప్రకారం, సంకేతం~$e$ గల ప్రమేయపు గణన ఇన్‌పుట్~$x$కు ఒక
విలువను ఉత్పత్తి చేస్తే, అప్పుడే మాత్రమే సార్వత్రిక పాక్షిక గణనీయ
ప్రమేయం~$\fn{Un}(e,x)$ నిర్వచితమవుతుంది. ఇదే పరిస్థితి వస్తుందో లేదో
మనం నిర్ణయించగలమా అని అడగడం సహజం. నిజానికి నిర్ణయించలేం. ట్యూరింగ్
యంత్ర గణనా నమూనాలో, ఇచ్చిన ట్యూరింగ్ యంత్రం ఇచ్చిన ఇన్‌పుట్‌పై ఆగుతుందో
లేదో గణనాత్మకంగా నిర్ణయించలేమని దీని అర్థం. అందువల్ల కింది సిద్ధాంతాన్ని
``నిలుపు సమస్య యొక్క అనిర్ణయనీయత'' అంటారు. కింద రెండు నిరూపణలు
ఇస్తాం. మొదటిది మన ముందు చర్చను కొనసాగిస్తుంది; రెండవది మరింత నేరుగా
ఉంటుంది.

\begin{thm}
\ollabel{thm:halting-problem}
\[
h(e,x)=
\begin{cases}
1 & \text{$\fn{Un}(e,x)$ నిర్వచితమైతే} \\
0 & \text{ఇతర సందర్భాల్లో}
\end{cases}
\]
అని ఉంచండి. అప్పుడు $h$ గణనీయం కాదు.
\end{thm}

\begin{proof}
$h$ గణనీయం అనుకుందాం. దాని వల్ల ఒక సార్వత్రిక గణనీయ ప్రమేయాన్ని
నిర్వచించవచ్చని చూపుతాం. ఇలా నిర్వచించండి:
\[
\fn{Un'}(e,x)=
\begin{cases}
\fn{Un}(e,x) & \text{$h(e,x)=1$ అయితే} \\
0 & \text{ఇతర సందర్భాల్లో.}
\end{cases}
\]
ఇప్పుడు $\fn{Un'}(e,x)$ సర్వనిర్వచిత ప్రమేయం; $h$ గణనీయమైతే ఇది
కూడా గణనీయం. ఉదాహరణకు, ఆదిమ పునరావృత్తిని వాడి $g$ను ఇలా
నిర్వచించవచ్చు:
\begin{align*}
g(0,e,x) & \simeq 0 \\
g(y+1,e,x) & \simeq \fn{Un}(e,x);
\end{align*}
అప్పుడు
\[
\fn{Un'}(e,x)\simeq g(h(e,x),e,x).
\]
$\fn{Un}(e,x)$ నిర్వచితమైన చోటల్లా $\fn{Un'}(e,x)$ దానితో
ఏకీభవిస్తుంది. కనుక పాక్షిక గణనీయ ప్రమేయాల్లో సర్వనిర్వచితంగా ఉండే
వాటికి $\fn{Un'}$ సార్వత్రికం. కానీ ఇది \olref[nou]{thm:no-univ}కు
వైరుధ్యం.
\end{proof}

\begin{proof}
$h(e,x)$ గణనీయం అనుకుందాం. $g$ను ఇలా నిర్వచించండి:
\[
g(x)=
\begin{cases}
  0 & \text{$h(x,x)=0$ అయితే} \\
  \fundefined & \text{ఇతర సందర్భాల్లో.}
\end{cases}
\]
$g$ పాక్షిక గణనీయ ప్రమేయం. ఉదాహరణకు, దాన్ని
$\umin{y}{h(x,x)=0}$గా నిర్వచించవచ్చు. కాబట్టి ఏదో~$e$కు ప్రతి~$x$కూ
$g(x)\simeq\fn{Un}(e,x)$ అవుతుంది. $e$ వద్ద $g$ నిర్వచితమా? నిర్వచితమైతే,
$g$ నిర్వచనం ప్రకారం $h(e,e)=0$ ($g$ నిర్వచితమవడానికి $h(e,e)$ విలువ
$0$ కావాలి). $h$ నిర్వచనం ప్రకారం, $\fn{Un}(e,e)$ అనిర్వచితమని దీని
అర్థం. ప్రతి~$x$కూ $g(x)\simeq\fn{Un}(e,x)$ అనే మన ఊహ ప్రకారం,
$g(e)$ అనిర్వచితం అవుతుంది; ఇది వైరుధ్యం. మరోవైపు $g(e)$ అనిర్వచితమైతే,
$h(e,e)\neq0$; అందువల్ల $h(e,e)=1$. దాని నుంచి $\fn{Un}(e,e)$
నిర్వచితమని వస్తుంది. కానీ $g(x)\simeq\fn{Un}(e,x)$ కనుక $g(e)$
కూడా నిర్వచితమవుతుంది. మళ్లీ వైరుధ్యం.
\sourcecorrection{OLTECOMTHY-003}{రెండవ నిరూపణలోని ఆంగ్ల మూలం, $g(e)$ నిర్వచితమైతే వచ్చే కారణాన్ని వివరిస్తూ, $h$ నిర్వచితమైతేనే అది $0$ విలువ తీసుకుంటుందని అసంగతంగా చెప్పింది. ఇక్కడ $h$ను గణనీయ ప్రమేయంగా ఊహించినందున అది సర్వనిర్వచితం; $g$ నిర్వచనం ప్రకారం $g(e)$ నిర్వచితమవడానికి $h(e,e)=0$ కావాలనే ఉద్దేశాన్ని స్పష్టం చేశాం.}
\end{proof}

\begin{tagblock}{TMs}
\begin{explain}
ఈ వాదనను ట్యూరింగ్ యంత్రాల పరంగా వర్ణించవచ్చు. ఒక ట్యూరింగ్ యంత్రం
$E$ యొక్క వివరణను, ఇన్‌పుట్~$x$ను తీసుకుని, $E$ ఇన్‌పుట్~$x$పై
ఆగుతుందో లేదో నిర్ణయించే ట్యూరింగ్ యంత్రం~$H$ ఉందనుకుందాం. అప్పుడు
ఒకే ఇన్‌పుట్~$x$ను తీసుకుని, సూచిక~$x$ గల యంత్రం $M_x$ ఇన్‌పుట్~$x$పై
ఆగుతుందో లేదో నిర్ణయించడానికి $H$ను నడిపి, దానికి విరుద్ధంగా చేసే
మరొక ట్యూరింగ్ యంత్రం~$G$ను నిర్మించవచ్చు. అంటే, ఇన్‌పుట్~$x$పై
$M_x$ ఆగుతుందని $H$ నివేదిస్తే $G$ అనంత చక్రంలోకి వెళుతుంది;
ఇన్‌పుట్~$x$పై $M_x$ ఆగదని $H$ నివేదిస్తే $G$ వెంటనే ఆగుతుంది.
తన స్వంత సూచికను ఇన్‌పుట్‌గా ఇచ్చినప్పుడు $G$ ఆగుతుందా? అది ఆగితేనే
ఆగదని పై వాదన చూపుతుంది---ఇది వైరుధ్యం. కనుక అటువంటి ట్యూరింగ్
యంత్రం~$H$ ఉందనే మన ఊహ తప్పు.
\end{explain}
\end{tagblock}

\end{document}
